\subsection*{“正在下降”与“下降得越来越慢”可以同时成立}
看位置函数的一阶导数，是在问速度；再对速度求导，则是在问加速度。两者的符号回答不同问题。一辆向左运动但正在减速的车，速度为负，加速度却可能为正。函数图像的下降与向上弯曲也可以同时发生。

以 $x^2$ 在负半轴为例，斜率 $2x$ 为负，但随着 $x$ 增加，它越来越接近零。图像仍下降，却下降得越来越慢，二阶导数为正。把“正二阶导数”误读为“函数递增”，就是把变化率与变化率的变化混在一起。

\subsection*{高阶信息用于更细的局部预测}
只知道一点的值，可以给常数预测；知道一阶导数，可以加直线修正；知道二阶导数，就能预测斜率如何偏离初值，从而加上二次项。泰勒近似会把这套逐层修正系统化。

对 $x^3$，二阶导数 $6x$ 在零点两侧变号，所以凹凸改变。对 $x^4$，二阶导数 $12x^2$ 在零点为零，却不变号，所以“二阶导数等于零”并不足以判断拐点。需要检查邻域行为，而非只代一个点。

试着画函数、斜率函数和二阶导函数三幅图，并用同一竖线标出一个输入。每幅图的纵坐标分别回答“在哪里、往哪变、变化率怎么变”。把它们分开阅读，才不会在更复杂的近似或动力系统中混淆层次。
